\documentclass{beamer}
\usepackage[utf8]{inputenc}
\usepackage[T1]{fontenc}
\usepackage[french]{babel}
\usepackage{amsmath}
\usepackage{graphicx}
\usepackage{booktabs}
\usetheme{Madrid}
\usecolortheme{default}

% Couleurs
\definecolor{titleblue}{RGB}{0, 51, 102}
\setbeamercolor{title}{fg=titleblue}
\setbeamercolor{subtitle}{fg=black}

\title{Une approche d'explicabilité des modèles d'apprentissage profond}
\subtitle{pour le contrôle d'accès aux données et aux traitements}
\author{TSAFACK NTEUDEM ERICK}
% \institute[Université de Dschang]{}
\date{\today}

\begin{document}
% ==================== SLIDE 1 : PAGE DE COUVERTURE ====================
\begin{frame}
\begin{center}
    \vspace{0.5cm}
    
    % Logo (optionnel)
    % \includegraphics[width=4cm]{logo_universite.png} 
    % \vspace{0.8cm}
    
    {\LARGE \textcolor{titleblue}{\textbf{RAPPORT D'AVANCEMENT}}} \\[0.2cm]
    {\large du Mémoire de Master} \\[0.2cm]
    
    {\Huge \textbf{Une approche d'explicabilité des modèles d'apprentissage profond pour le contrôle d'accès aux données et aux traitements}} \\[0.2cm]
    
    {\Large \textbf{Présenté par :}} \\[0.3cm]
    {\Large \insertauthor} \\[1cm]
    
    {\large Encadré par : Dr. AZANGUEZET Martin} \\[1.2cm]
    
    {\large \insertdate \quad\quad Version 1.0}
    
\end{center}
\end{frame}

\begin{frame}
\frametitle{1. Introduction et Motivation (Toward Deep Learning Based Access Control)}
\framesubtitle{Problèmes des approches classiques de contrôle d’accès au donnée}

\begin{itemize}
    \item Les modèles classiques (RBAC, ABAC, ReBAC) nécessitent :
    \begin{itemize}
        \item \textbf{Attribute Engineering} : conception manuelle d'attributs pertinents à partir de métadonnées brutes
        \item \textbf{Policy Engineering} : création manuelle ou semi-automatique de règles d'accès
    \end{itemize}
    \item Limitations majeures :
    \begin{itemize}
        \item Coût élevé et erreurs humaines (over/under-provisioning)
        \item Mauvaise généralisation sur des utilisateurs/ressources non vus
        \item Difficile à maintenir dans les systèmes dynamiques, à grande échelle (Cloud, IoT, etc.)
    \end{itemize}
    \item \textbf{Proposition} : \textbf{Deep Learning Based Access Control (DLBAC)}
\end{itemize}

\begin{center}
\textbf{Idée clé} : Remplacer les règles par un réseau de neurones entraîné directement sur les \textbf{métadonnées brutes}.
\end{center}

\end{frame}

\begin{frame}
\frametitle{2. Le Concept DLBAC}
\framesubtitle{Approche révolutionnaire}

\begin{itemize}
    \item \textbf{Entrée} : Métadonnées brutes des utilisateurs et ressources (join\_date, department, logs, etc.)
    \item \textbf{Modèle} : Réseau de neurones profond (indépendant de l'architecture)
    \item \textbf{Sortie} : Décision d'accès (grant/deny) pour chaque opération
\end{itemize}

% \begin{center}
% \includegraphics[width=0.9\textwidth]{figures/comparison.png} % Remplacer par vraie image si disponible
% \end{center}

\textbf{Avantages principaux :}
\begin{itemize}
    \item Plus besoin d'ingénierie d'attributs ni de règles d'accès
    \item Meilleure généralisation grâce au Deep Learning
    \item Approche \textbf{end-to-end} et dynamique
\end{itemize}

\textbf{DLBAC$_\alpha$} : Prototype proposé et évalué dans le papier.

\end{frame}

\begin{frame}
\frametitle{3. Évaluation et Résultats}
\framesubtitle{DLBAC$_\alpha$ vs Approches existantes}

\begin{itemize}
    \item \textbf{Datasets} : 2 réels (Amazon) + 8 synthétiques de complexité croissante
    \item \textbf{Modèles testés} :
    \begin{itemize}
        \item DLBAC$_\alpha$ (ResNet, DenseNet, Xception)
        \item ML classiques (RF, SVM, MLP)
        \item Policy Mining (XuStoller, Rhapsody, EPDE-ML)
    \end{itemize}
\end{itemize}

\textbf{Résultats clés :}
\begin{itemize}
    \item DLBAC$_\alpha$ obtient les meilleurs scores \textbf{F1}, TPR et généralisation
    \item Meilleur équilibre entre \textit{over-provisioning} et \textit{under-provisioning}
    \item Plus robuste face à la complexité des données
\end{itemize}

\begin{center}
\small Meilleure précision + généralisation que le policy mining et le ML classique
\end{center}

\end{frame}

\begin{frame}
\frametitle{4. Explicabilité, Perspectives \& Conclusion}
\framesubtitle{Comprendre les décisions du modèle}

\begin{itemize}
    \item \textbf{Explainability} :
    \begin{itemize}
        \item \textit{Integrated Gradients} : importance de chaque métadonnée (local + global)
        \item \textit{Knowledge Transferring} : extraction d'arbres de décision approximatifs
    \end{itemize}
    \item Permet de justifier les décisions et de modifier des métadonnées pour ajuster les accès
\end{itemize}

\vspace{0.3cm}
\textbf{Directions futures :}
\begin{itemize}
    \item Explicabilité Post-hoc
    \item Administration (fine-tuning, life-long learning)
    \item Robustesse (attaques adversariales, biais)
    \item Intégration hybride avec RBAC/ABAC
\end{itemize}

\begin{center}
\textbf{Conclusion :} DLBAC offre une approche prometteuse, automatisée et performante pour l'Access Control du futur.
\end{center}

\end{frame}

\begin{frame}
\frametitle{Pourquoi les explications Post-Hoc sont insuffisantes pour DLBAC ? (A Comprehensive Survey on Self-Interpretable
Neural Networks)}

\begin{itemize}
    \item DLBAC est un système de \textbf{décision critique} en sécurité
    \item Les méthodes post-hoc sont :
    \begin{itemize}
        \item \alert{Fragiles} et peu robustes face aux attaques
        \item \alert{Peu fidèles} au véritable fonctionnement du réseau
        \item \alert{Instables} : explications variables pour des entrées similaires
        \item Computationnellement coûteuses
    \end{itemize}
    \item Risque majeur : fausse confiance dans les décisions d’accès
\end{itemize}

\begin{center}
\textbf{\Large Nécessité d’une architecture Self-Interpretable}
\end{center}

\end{frame}

\begin{frame}
\frametitle{Choix des Explications pour DLBAC (AI Explainability 360: Impact and Design)}
\framesubtitle{Contrôle d'accès basé sur Deep Learning}
\begin{itemize}Quel type d'explication choisir ?
    \item des Contrefactuelles
    \item des règles de décisions
\end{itemize}

\begin{block}{Pourquoi les explications \textbf{Contrefactuelles} ?}
\begin{itemize}
    \item Répondent à la question : \textbf{« Que faut-il changer pour obtenir/refuser l'accès ? »}
    \item Très \textbf{actionnables} pour les utilisateurs et administrateurs
    \item Permettent de comprendre les \textbf{frontières de décision} du modèle
    \item Essentielles en sécurité : expliquer un refus d'accès de manière constructive
    \item Exemple : « Ajoutez le rôle X ou changez le département Y pour obtenir l'accès »
\end{itemize}
\end{block}
\end{frame}

\begin{frame}
\frametitle{Choix des Explications pour DLBAC (AI Explainability 360: Impact and Design)}
\begin{block}{Pourquoi les explications \textbf{basées sur des Règles} ?}
\begin{itemize}
    \item Fournissent une \textbf{logique globale et auditable} du modèle
    \item Faciles à comprendre par les non-experts et les régulateurs
    \item Permettent la \textbf{vérification formelle} des politiques d'accès
    \item Compatibles avec les exigences de conformité (RGPD, audits de sécurité)
    \item Exemple : « Accès accordé si (département = IT) ET (niveau senior) »
\end{itemize}
\end{block}

\begin{center}
\textbf{Combinaison idéale pour DLBAC :} \\
Règles (globale + auditable) + Contrefactuelles (locale + actionnable)
\end{center}

\end{frame}

\begin{frame}
\frametitle{HyConEx : Hypernetwork Classifier avec Explications Contrefactuelles}

\begin{itemize}
    \item Modèle Deep Learning conçu pour les \textbf{données tabulaires}
    \item Combine dans une seule architecture :
    \begin{itemize}
        \item Prédiction de classe performante
        \item Génération de \textbf{contrefactuels plausibles} pour toutes les classes alternatives
    \end{itemize}
    \item Tout est produit en \textbf{un seul forward pass} (très efficace)
\end{itemize}

\begin{center}
\textbf{Idée clé :} Un modèle qui prédit \textbf{et} explique simultanément comment modifier l'entrée pour changer la décision.
\end{center}

\end{frame}

\begin{frame}
\frametitle{Positionnement de HyConEx}

\begin{itemize}
    \item Modèles interprétables classiques (arbres, règles) : bonne explicabilité mais performance limitée
    \item Modèles DL + explications post-hoc (LIME, SHAP, CEM, PPCEF) : haute performance mais explications lentes et externes
    \item \textbf{HyConEx} : première approche qui intègre \textbf{prédiction + explication contrefactuelle} dans un seul modèle
\end{itemize}

\begin{center}
\textbf{Avantage majeur pour le contrôle d'accès :} \\
Explication actionnable immédiatement (« que faut-il changer pour obtenir/refuser l'accès ? »)
\end{center}

\end{frame}

\begin{frame}
\frametitle{Méthodologie de HyConEx}

\begin{itemize}
    \item \textbf{Hypernetwork} $H(x;\theta)$ : génère pour chaque entrée $x$ une matrice de poids $W$
    \item $W$ définit un \textbf{classifieur linéaire local} → frontière de décision locale
    \item Les lignes de $W$ servent à générer des contrefactuels par translation : $x' = x - W_m$
\end{itemize}

\vspace{0.4cm}

\begin{block}{Fonction de perte globale}
\begin{itemize}
    \item Perte de classification de l'exemple d'entrée
    \item Pour chaque classe alternative :
    \begin{itemize}
        \item Perte de classification du contrefactuel
        \item Perte de proximité ($\|x - x'\|$)
        \item Perte de plausibilité (via Normalizing Flow conditionnel)
    \end{itemize}
\end{itemize}
\end{block}

\begin{center}
\textbf{Résultat :} Prédictions + Contrefactuels plausibles et actionnables en un seul passage
\end{center}

\end{frame}

\begin{frame}
\frametitle{HyperLogic : Amélioration de l'apprentissage de règles avec HyperNetworks}

\begin{itemize}
    \item Approche basée sur les \textbf{hypernetworks} pour générer les poids d'un réseau principal d'apprentissage de règles
    \item Objectif : Augmenter la \textbf{diversité} et la \textbf{précision} des règles tout en préservant l'interprétabilité
    \item Permet de générer \textbf{plusieurs ensembles de règles} de qualité en une seule session d'entraînement
    \item Compatible avec les méthodes différentiables d'apprentissage de règles existantes
\end{itemize}

\begin{center}
\textbf{Idée clé :} Utiliser la flexibilité des hypernetworks pour contourner les limitations structurelles des réseaux de règles
\end{center}

\end{frame}

\begin{frame}
\frametitle{Positionnement de HyperLogic}

\begin{itemize}
    \item Les réseaux de règles classiques sont souvent trop rigides et peu expressifs
    \item Les approches différentiables actuelles manquent de diversité dans les règles apprises
    \item \textbf{HyperLogic} apporte :
    \begin{itemize}
        \item Plus grande flexibilité grâce à la génération dynamique des poids
        \item Meilleure capture des patterns hétérogènes dans les données
        \item Plusieurs ensembles de règles candidats en un seul entraînement
    \end{itemize}
\end{itemize}

\begin{center}
\textbf{Avantage majeur :} Meilleure expressivité sans sacrifier l'interprétabilité
\end{center}

\end{frame}

\begin{frame}
\frametitle{Méthodologie de HyperLogic}

\begin{itemize}
    \item \textbf{Hypernetwork} : génère les poids $\theta$ du réseau principal à partir d'un bruit gaussien
    \item \textbf{Réseau principal} : réseau de règles différentiable (inspiré de DR-Net)
        \begin{itemize}
            \item Couche "Rule" : apprend des conjonctions logiques
            \item Couche "OR" : combine les règles
        \end{itemize}
    \item Deux stratégies de combinaison des poids possibles (fusion pondérée ou génération partielle)
    \item Régularisation : sparsité + diversité (KL divergence)
\end{itemize}

\begin{center}
\textbf{Résultat :} Apprentissage de règles plus diverses, précises et concises
\end{center}

\end{frame}

\begin{frame}
\frametitle{Méthodologie de HyperLogic}

\begin{itemize}
    \item \textbf{Hypernetwork} $G(\epsilon)$ génère les poids $\theta = (w, u)$ du réseau principal :
    $$
    \theta = G(\epsilon), \quad \epsilon \sim \mathcal{N}(0, I)
    $$
    
    \item \textbf{Réseau principal (inspiré de DR-Net)} :
    \begin{itemize}
        \item \textbf{Couche Rule} apprend des conjonctions logiques : pour chaque règle $k$
        $$
        o_k = h\left( \sum_{d=1}^D w_{d,k} x_d - \sum_{d=1}^D |w_{d,k}| \right)
        $$
        avec $h(u) = \exp\left(-\frac{u^2}{\tau^2}\right)$ (approximation différentiable)
        
        \item \textbf{Couche OR} : Combine les règles
        $$
        f(x) = \sum_{k=1}^K u_k \, o_k
        $$
    \end{itemize}
    
    \item Modèle final :
    $$
    f_\mu(x) = \mathbb{E}_{\theta \sim \mu} \left[ \sum_{k=1}^K u_k \, h(\dots) \right] \approx \frac{1}{M} \sum_{m=1}^M f_{\theta^m}(x)
    $$
\end{itemize}

\end{frame}

\begin{frame}
\frametitle{Méthodologie de HyperLogic}

\begin{itemize}
    \item \textbf{1. Stratégies de combinaison des poids}
    
    \begin{block}{Fusion pondérée (stratégie principale)}
    $$
    W_{\text{final}} = \alpha \cdot W_{\text{hyper}} + (1 - \alpha) \cdot W_{\text{main}}
    $$
    où $\alpha \in [0,1]$ est un paramètre \textbf{apprenable}
    \end{block}
    
    \vspace{0.2cm}
    
    \begin{block}{Génération partielle}
    La hypernetwork ne génère que certains modules du réseau principal (ex. : seulement la Rule Layer) pour les très hautes dimensions.
    \end{block}

    \item \textbf{2. Régularisation}
    $$
    \ell_{\text{reg}} = \lambda_1 \, D_{\text{KL}}(\mu \, || \, \mu_0) + \lambda_2 \, \mathbb{E}_{\mu}\left[ \sum_{k=1}^K |u_k| \right]
    $$
    
    \begin{itemize}
        \item $\lambda_1 D_{\text{KL}}$ $\rightarrow$ \textbf{encourage la diversité} des ensembles de règles
        \item $\lambda_2 \mathbb{E}[|u_k|]$ $\rightarrow$ \textbf{encourage la sparsité} (règles plus concises)
    \end{itemize}
\end{itemize}


\begin{center}
\textbf{Conséquence :} Meilleure expressivité tout en maintenant l'interprétabilité
\end{center}

\end{frame}

\begin{frame}{Méthodologie de HyperLogic}
    \begin{block}{Fonction de perte}
$\ell_{\text{task}} + \lambda_1 \, D_{\text{KL}}(\mu \| \mu_0) + \lambda_2 \, \mathbb{E}_\mu \left[ \sum |u_k| \right]$
\end{block}
\end{frame}

\begin{frame}
\frametitle{Perspectives pour mon modèle DLBAC (Règles + Contrefactuelles)}

\begin{itemize}
    \item \textbf{HyperLogic} excelle dans l'apprentissage de \textbf{règles interprétables} et diversifiées
    \item \textbf{HyConEx} excelle dans la génération rapide de \textbf{contrefactuels plausibles}
    \item \textbf{Idée d'intégration pour DLBAC :}
    \begin{itemize}
        \item Utiliser un hypernetwork pour générer plusieurs ensembles de règles (globale + auditable)
        \item Combiner avec une branche contrefactuelle inspirée de HyConEx (locale + actionnable)
        \item Un seul modèle produisant à la fois :
        \begin{itemize}
            \item Règles logiques pour l’audit et la conformité
            \item Contrefactuels pour expliquer les décisions d’accès
        \end{itemize}
    \end{itemize}
\end{itemize}

\begin{center}
\textbf{Objectif final :} Un modèle DLBAC, performant, interprétable et actionnable
\end{center}

\end{frame}

% ==================== Modèle unifié (règles + contrefactuels) : 5 slides théoriques ====================

\begin{frame}
\frametitle{Un modèle unifié pour DLBAC}
\framesubtitle{Règles globales et contrefactuels locaux}

\begin{itemize}
    \item On cherche une fonction $f$ sur des données \textbf{tabularies} telle que, pour chaque instance $x$ :
    \begin{itemize}
        \item $f(x)$ encode une \textbf{décision de classe} (accès / refus, ou multi-classes) ;
        \item une \textbf{lecture en règles} révèle des conjonctions de littéraux sur une représentation discrète ;
        \item une \textbf{branche contrefactuelle} propose des modifications minimales de $x$ pour atteindre une classe cible $c \neq \hat{y}(x)$.
    \end{itemize}
    \item \textbf{Hypernetwork} $H$ : à partir d'une représentation latente dérivée de l'instance (éventuellement agrégée sur plusieurs entrées), $H$ produit deux vecteurs de paramètres
    \[
    (\theta_{\text{rule}}, \theta_{\text{cf}}) = H(\cdot)\,,
    \]
    qui spécialisent respectivement le classifieur à base de règles et le générateur de contrefactuels.
\end{itemize}

\begin{center}
\textbf{Synthèse} : une seule passe avant combine \textbf{politique locale} (paramètres instantanés) et \textbf{explication actionnable}.
\end{center}

\end{frame}

\begin{frame}
\frametitle{Représentation discrète bipolarisée}
\framesubtitle{De $\mathbb{R}^p$ à $\{-1,+1\}^D$}

\begin{itemize}
    \item Chaque coordonnée continue est discrétisée en $B$ intervalles (par exemple via quantiles empiriques), ce qui définit $B$ \textbf{littéraux exclusifs} par variable.
    \item L'instance est plongée dans un espace $\{-1,+1\}^D$ où $D$ est la somme des nombres d'intervalles par attribut : exactement un littéral par variable actif à $+1$, les autres à $-1$.
    \item Ce plongement aligne la géométrie de l'entrée sur la \textbf{couche Rule} de type HyperLogic / DR-Net, où les poids $w_k$ agissent comme des conjonctions relaxées sur ces littéraux.
    \item Pour mesurer la \textbf{proximité} d'un contrefactuel $x'$ vis-à-vis de $x$, on peut repasser par le domaine continu (barycentre des intervalles actifs) et utiliser une norme $\|\cdot\|_1$ ou $\ell_2$.
\end{itemize}

\begin{block}{Interprétation DLBAC}
Chaque dimension de $\{-1,+1\}^D$ est un \textbf{attribut discretisé} (niveau de séniorité, tranche d'ancienneté, etc.), lisible pour un auditeur.
\end{block}

\end{frame}

\begin{frame}
\frametitle{Classifieur à règles induit par $\theta_{\text{rule}}$}
\framesubtitle{Même structure que DR-Net / HyperLogic}

\begin{itemize}
    \item Le vecteur $\theta_{\text{rule}}$ se décompose en poids de règles $w_k \in \mathbb{R}^D$, biais $b_k$, matrice de sortie $U \in \mathbb{R}^{K \times C}$ et biais de classe.
    \item Score de la $k$-ième règle (activation douce) :
    \[
    u_k(x) = x^\top w_k - \|w_k\|_1 + b_k\,, \qquad
    o_k(x) = \exp\!\left(-\frac{u_k(x)^2}{\tau^2}\right), \quad \tau > 0\,.
    \]
    \item Logits multi-classes :
    \[
    z(x) = O(x)\, U + b_{\text{out}}\,, \qquad O(x) = \bigl(o_1(x),\ldots,o_K(x)\bigr)\,.
    \]
    \item Les $o_k$ jouent le rôle d'une \textbf{couche OR} sur des conjonctions : $u_k$ pénalise les littéraux non alignés avec les signes de $w_k$.
\end{itemize}

\begin{center}
\textbf{Lecture} : les colonnes $w_k$ mettent en évidence quels littéraux $\pm 1$ soutiennent chaque règle ; $U$ route chaque règle vers les classes.
\end{center}

\end{frame}

\begin{frame}
\frametitle{Contrefactuels dans $\{-1,+1\}^D$}
\framesubtitle{Réseau conditionné par la classe cible}

\begin{itemize}
    \item On fixe une classe désirée $c$ et on forme le vecteur one-hot $e_c \in \mathbb{R}^C$.
    \item Un réseau auxiliaire $g_{\theta_{\text{cf}}}$ (paramétré par l'hypernetwork) lit la paire $(x, e_c)$ et produit, pour chaque dimension $j$, un logit $\ell_j$ puis $p_j = \sigma(\ell_j) \approx \mathbb{P}(x'_j = +1)$.
    \item On définit un contrefactuel \textbf{dur} $x' \in \{-1,+1\}^D$ par seuillage, tout en rétro-propageant à travers une relaxation continue (estimateur straight-through) pour optimiser les flips.
    \item Pour chaque $c$, on impose que $f(x')$ classe $c$ (entropie croisée sur $z(x')$) et que $x'$ reste \textbf{proche} de $x$ (pénalité sur le nombre ou le coût des changements de littéraux).
\end{itemize}

\begin{block}{Lien avec HyConEx}
Même esprit : \textbf{une décision} et \textbf{un chemin de modification} co-entraînés, ici dans un espace discret compatible avec les règles.
\end{block}

\end{frame}

\begin{frame}
\frametitle{Objectif d'apprentissage}
\framesubtitle{Perte composite}

\begin{itemize}
    \item Perte totale (scalaires $\lambda_{\text{cf}}, \lambda_{\text{flip}}, \lambda_{\text{sparse}} \geq 0$) :
\end{itemize}

\vspace{-0.2cm}
\[
\begin{aligned}
\mathcal{L}
    &= \underbrace{\mathrm{CE}\bigl(z(x), y\bigr)}_{\text{classification}} \\[0.35em]
    &\quad + \lambda_{\text{cf}} \sum_{c \neq y} \mathrm{CE}\bigl(z(x'_c), c\bigr) \\[0.35em]
    &\quad + \lambda_{\text{flip}}\, \mathbb{E}\bigl[d(x, x'_c)\bigr] \\[0.35em]
    &\quad + \lambda_{\text{sparse}}\, \frac{1}{KD}\sum_{k,j} |w_{k,j}|\,.
\end{aligned}
\]

\vspace{0.2cm}
\begin{itemize}
    \item $d$ : coût des modifications (nombre de littéraux changés, ou distance en espace continu reconstruit).
    \item Terme sparse : favorise des règles \textbf{courtes}, plus simples à auditer.
\end{itemize}

\end{frame}

\begin{frame}
\frametitle{Évaluation et lecture des politiques}
\framesubtitle{Contrefactuels et règles}

\begin{itemize}
    \item \textbf{Qualité des contrefactuels} : taux de validité $\mathbb{P}(\arg\max z(x'_c) = c)$, parcimonie moyenne $\mathbb{E}[\|x - x'\|_0]$, proximité $\mathbb{E}[\|x - x'\|_1]$ après reconstruction.
    \item \textbf{Lecture des règles} : pour chaque $k$, extraire les indices $j$ avec $|w_{k,j}|$ notable et lire le littéral actif ($+1$ ou $-1$) comme condition ; $U$ indique la classe dominante de la règle.
\end{itemize}

\vspace{0.4cm}
\begin{center}
\textbf{But DLBAC} : précision élevée, règles inspectables, contrefactuels \textbf{plausibles} et \textbf{peu coûteux} pour l'utilisateur.
\end{center}

\end{frame}

\end{document}